\documentclass[10pt,a4paper]{article}
\usepackage{amsmath,amssymb,amsthm,graphicx}
\usepackage{hyperref}
\usepackage{latexsym}
\usepackage{amssymb}
\usepackage{amsthm}
\usepackage{ulem}
\usepackage{enumitem}
\setlist{itemsep=0.4pt,topsep=0pt}
\parindent 0pt
\parindent 0pt
\parskip 12pt
\newcommand{\HCF}{\mathrm{hcf}}
\newcommand{\ZZ}{\mathbb{Z}}
\newcommand{\rk}{\mathrm{rk}}
\DeclareMathOperator{\len}{length}
\DeclareMathOperator{\bl}{b}
\DeclareMathOperator{\cl}{cl}
\DeclareMathOperator{\di}{d}
\DeclareMathOperator{\Cl}{Cl}
\DeclareMathOperator{\acl}{acl}
\DeclareMathOperator{\tp}{tp}
\DeclareMathOperator{\Th}{Th}
\newcommand{\Cor}{\mathcal{C}_\mathrm{or}}
\newcommand{\Cun}{\mathcal{C}_\mathrm{un}}
\newcommand{\mons}{\mathbb{M}}
\newcommand{\strong}{\mathrm{strong}}
\newcommand{\axioms}{\Delta_{\mathrm{FSTS}}}
\newtheoremstyle{note}% hnamei
{8pt}% hSpace abovei
{14pt}% hSpace belowi
{\normalfont}% hBody fonti
{}% hIndent amounti1
{\bfseries}% hTheorem head fonti
{ }% hPunctuation after theorem headi
{.5em}% hSpace after theorem headi2
{}% hTheorem head spec (can be left empty, meaning ‘normal’)i
\theoremstyle{note}
\newtheorem{thm}{Theorem}[section]
%\newtheorem{theorem}{Theorem}
\newtheorem{lemma}[thm]{Lemma} 
\newtheorem{df}[thm]{Definition}
\newtheorem{rmk}[thm]{Remark} 
\newtheorem{prop}[thm]{Proposition}
\newtheorem{fact}[thm]{Fact}
\newtheorem{cor}[thm]{Corollary}
\newtheorem{conj}[thm]{Conjecture}
\newtheorem{ex}[thm]{Example}
\newtheorem{qu}[thm]{Question}
\newtheorem*{notation}{Notation} 
\title{An HF-ordering is over the set it closes}
\author{Silvia Barbina, Enrique Casanovas}
\date{arxiv:2411.13723v3; Creative Commons Attribution 4.0 International}
\begin{document}
\maketitle
\noindent\textit{Editorial note.} The counted statement is the author's labelled environment at source0.tex lines 346-347 and its proof environment opens at line 348; the pairing is by adjacency and the retained proof range is the whole environment.\par\medskip
\noindent\textit{Editorial note.} This source is newly approved for this line; the preamble is built from the author's own package lines, macro block and theorem-environment declarations (source lines 4-65), with the author's optioned amsmath line omitted so that it cannot clash with the wrapper's amsmath.\par\medskip
\noindent\textit{Editorial note.} The source carries no label-number dump and the retained counter values are left empty, so the printed numbering of the card is the wrapper's default; the identity of the counted item is its author label.\par\medskip
\noindent\textit{Editorial note.} Metadata provenance: the source record of the pinned copy and the licence evidence file are both present; the licence is CC BY 4.0 with status supported.\par\medskip
\noindent\textit{Editorial note.} r002 note: the pre-publication gate compares the card's source\_version with the source\_version recorded in the licence evidence of the same source. This source's evidence file records exactly "2411.13723v3", while the card had carried "arXiv:2411.13723v3 (CC BY 4.0)" copied from the source record; r002 aligns the card with the evidence and changes nothing else. The evidence file itself is not modified.\par\medskip
\noindent\textit{Editorial scope.} One original result of the pinned author manuscript, counted once: an HF-ordering is over the set it closes (author label Feb\_23\_2024\_7).\par\medskip
\noindent\textit{Allowed standard prerequisites (editorial).}\begin{itemize}
\item HF-ordering over a subset. An ordering is over A when every element outside A has at most one pair of smaller elements forming a block with it.. It is the property the lemma establishes.
\item Closure of a subset. The closure of A under the ordering is the smallest set containing A and closed under forming blocks whose other members are already present and smaller.. The hypothesis is that A equals its own closure.
\end{itemize}\medskip
\begin{lemma}\label{Feb_23_2024_7} If $<$ is an HF-ordering of $B$  and $A=\cl_<(A)\subseteq B$, then  $<$ is an HF-ordering of $B$ over $A$. 
\end{lemma}

\begin{proof}   If not, there is some $b\in B\smallsetminus A$ and there are distinct $b_1,b_2,b_3,b_4$, such that  $\{b,b_1,b_2\}$,  $\{b,b_3,b_4\}$  are  blocks and, for each $i$, either $b_i < b$ or $b_i \in A$.   Without loss of generality,  $b_2 <b_1$. We claim that $b_1 < b$.  If not, then   $b_1\in A$ and $b <b_1$,  and since $A=\cl_<(A)$, we get $b \in A$, a contradiction. Hence,  $b>b_1>b_2$. By a similar argument, $b> b_3,b_4$.  But then $<$ is not an HF-ordering.
\end{proof}

\end{document}
